\chapter{KIẾN THỨC CƠ SỞ VÀ CÁC CÔNG TRÌNH LIÊN QUAN}

Chương này trình bày các khái niệm nền tảng làm cơ sở cho nghiên cứu: mô hình ngôn ngữ nhỏ, kỹ thuật tinh chỉnh LoRA/QLoRA, đặc thù của NLP pháp lý tiếng Việt và tổng quan các công trình liên quan.

\section{Mô hình ngôn ngữ nhỏ}

Mô hình ngôn ngữ lớn (LLM) dựa trên kiến trúc Transformer đã đạt năng lực ấn tượng trên nhiều tác vụ, nhưng kích thước hàng chục đến hàng trăm tỷ tham số khiến chúng khó triển khai trong thực tế. Mô hình ngôn ngữ nhỏ (Small Language Model - SLM), thường có quy mô từ vài trăm triệu đến khoảng mười tỷ tham số, xuất hiện như một hướng đi cân bằng giữa năng lực và chi phí. Ưu điểm của SLM bao gồm chi phí triển khai thấp, khả năng triển khai cục bộ để bảo mật dữ liệu, và đặc biệt phù hợp để tinh chỉnh chuyên sâu cho một lĩnh vực hẹp như pháp lý.

Qwen3~\cite{yang2025qwen3} là họ mô hình ngôn ngữ mã nguồn mở do Alibaba phát triển, trải rộng từ 0.6B đến 235B tham số. Qwen3 hỗ trợ 119 ngôn ngữ (bao gồm tiếng Việt), tích hợp cả chế độ suy luận có tư duy (thinking mode) lẫn chế độ phản hồi nhanh và được phát hành dưới giấy phép Apache 2.0. Trong khuôn khổ cuộc thi VLSP2025-LegalSML, ban tổ chức cung cấp mô hình \texttt{qwen3-4b-legal-pretrain} -- một phiên bản Qwen3-4B đã được thích ứng miền qua tiền huấn luyện trên kho dữ liệu pháp lý tiếng Việt. Việc kế thừa mô hình này giúp nghiên cứu tập trung vào giai đoạn tinh chỉnh theo chỉ dẫn thay vì xây dựng tri thức miền từ đầu.

\section{Kỹ thuật tinh chỉnh LoRA và QLoRA}

Khi muốn dạy một mô hình ngôn ngữ lớn (LLM) thực hiện một tác vụ mới, cách truyền thống là \textbf{tinh chỉnh toàn phần} (Full Fine-Tuning): cập nhật tất cả các tham số của mô hình. Tuy nhiên, với mô hình có hàng tỷ tham số, cách này đòi hỏi bộ nhớ GPU rất lớn (thường trên 80GB) và tốn nhiều thời gian~\cite{hf_peft}. Để giải quyết vấn đề này, các kỹ thuật \textbf{tinh chỉnh hiệu quả tham số} (Parameter-Efficient Fine-Tuning - PEFT) ra đời, cho phép chỉ cập nhật một phần nhỏ tham số trong khi ``đóng băng'' phần còn lại~\cite{lightning_lora}.

\subsection{LoRA -- Low-Rank Adaptation}

LoRA (Low-Rank Adaptation)~\cite{hu2022lora} là một kỹ thuật PEFT phổ biến, dựa trên một ý tưởng đơn giản: khi tinh chỉnh mô hình cho một tác vụ cụ thể, lượng thay đổi cần thiết trong các trọng số thực ra có thể được biểu diễn trong một không gian có chiều thấp hơn nhiều so với không gian gốc~\cite{towardsdatascience_lora}.

\textbf{Ý tưởng chính.} Thay vì cập nhật trực tiếp ma trận trọng số $W$ có kích thước lớn ($d \times k$, với $d$ và $k$ có thể lên tới hàng nghìn), LoRA ``chèn'' thêm hai ma trận nhỏ $A$ và $B$ vào mạng. Lượng thay đổi $\Delta W$ được tính bằng tích của hai ma trận này:
\begin{equation}
W_{\text{mới}} = W_{\text{gốc}} + \Delta W = W_{\text{gốc}} + B \times A
\end{equation}
trong đó $A$ có kích thước $r \times k$, $B$ có kích thước $d \times r$, và $r$ (gọi là \textbf{rank}) là một số nhỏ, thường từ 4 đến 64~\cite{hf_lora_conceptual}.

\textbf{Tại sao tiết kiệm bộ nhớ?} Số tham số cần huấn luyện giảm đáng kể. Ví dụ với $d = k = 4096$ và $r = 16$:
\begin{itemize}
    \item Tinh chỉnh toàn phần: $4096 \times 4096 = 16.777.216$ tham số
    \item LoRA: $(4096 \times 16) + (16 \times 4096) = 131.072$ tham số
\end{itemize}
Như vậy LoRA giảm khoảng \textbf{128 lần} số tham số cần cập nhật~\cite{lightning_lora}.

\textbf{Khởi tạo.} Ma trận $A$ được khởi tạo ngẫu nhiên theo phân phối Gaussian, còn $B$ được khởi tạo bằng 0. Điều này đảm bảo rằng ban đầu $\Delta W = B \times A = 0$, tức mô hình bắt đầu với đúng hành vi của mô hình gốc~\cite{towardsdatascience_lora}.

\textbf{Hệ số tỉ lệ alpha.} LoRA sử dụng thêm một hệ số $\alpha$ để điều chỉnh mức độ ảnh hưởng của phần adapter:
\begin{equation}
h = W_{\text{gốc}} \cdot x + \frac{\alpha}{r} \cdot B \cdot A \cdot x
\end{equation}
Tỉ lệ $\frac{\alpha}{r}$ giúp kiểm soát biên độ thay đổi: $\alpha$ lớn hơn $r$ sẽ tăng ảnh hưởng của adapter, và ngược lại~\cite{hf_lora_conceptual}.

\textbf{Hợp nhất sau huấn luyện.} Sau khi huấn luyện xong, có thể ``gộp'' $\Delta W = BA$ vào trọng số gốc: $W \leftarrow W + BA$. Khi đó, mô hình hoạt động như một mô hình thông thường mà không có thêm chi phí tính toán khi suy luận~\cite{lightning_lora}.

\subsection{QLoRA -- Quantized LoRA}

QLoRA~\cite{dettmers2023qlora} là phiên bản nâng cấp của LoRA, cho phép tinh chỉnh các mô hình rất lớn trên GPU có bộ nhớ hạn chế bằng cách kết hợp LoRA với kỹ thuật \textbf{lượng tử hóa} (quantization)~\cite{hf_qlora}.

\textbf{Lượng tử hóa là gì?} Lượng tử hóa là kỹ thuật giảm độ chính xác của các số trong mô hình. Thay vì lưu mỗi tham số dưới dạng số thực 32-bit (float32) hoặc 16-bit (float16), ta có thể nén xuống 8-bit hoặc thậm chí 4-bit. Điều này giúp giảm đáng kể dung lượng bộ nhớ cần thiết~\cite{hf_quantization}.

\textbf{Cách QLoRA hoạt động.} QLoRA áp dụng ba kỹ thuật chính~\cite{dettmers2023qlora}:

\begin{enumerate}
    \item \textbf{Lượng tử hóa 4-bit NF4}: Mô hình gốc được nén xuống 4-bit bằng định dạng NormalFloat (NF4), được thiết kế riêng cho phân phối trọng số của mạng neural. Mô hình 4B tham số chỉ chiếm khoảng 2GB thay vì 8GB (ở float16)~\cite{hf_qlora}.

    \item \textbf{Adapter LoRA ở độ chính xác cao}: Trong khi mô hình gốc được giữ ở 4-bit và đóng băng, các ma trận adapter $A$ và $B$ vẫn được huấn luyện ở độ chính xác 16-bit (bfloat16 hoặc float16). Điều này đảm bảo quá trình học vẫn chính xác~\cite{lightning_lora}.

    \item \textbf{Giải nén khi cần}: Khi tính toán, trọng số 4-bit được tạm thời giải nén lên 16-bit để thực hiện phép nhân ma trận, sau đó gradient chỉ cập nhật vào adapter. Quá trình này diễn ra tự động và không ảnh hưởng đến kết quả~\cite{hf_qlora}.
\end{enumerate}

\textbf{Lợi ích thực tế.} Nhờ QLoRA, việc tinh chỉnh mô hình 4 tỷ tham số có thể thực hiện trên GPU Tesla T4 với chỉ 16GB VRAM -- loại GPU miễn phí trên các nền tảng như Google Colab hay Kaggle. Đây là nền tảng kỹ thuật cho phép toàn bộ quy trình của nghiên cứu này chạy được trên môi trường notebook miễn phí~\cite{hf_peft}.

\section{Tinh chỉnh theo chỉ dẫn}

Tinh chỉnh theo chỉ dẫn (Instruction Tuning) là giai đoạn huấn luyện sau tiền huấn luyện, trong đó mô hình được tinh chỉnh trên các cặp (chỉ dẫn, phản hồi mong muốn) để học cách tuân theo hướng dẫn của người dùng. Dữ liệu được định dạng theo mẫu hội thoại chuẩn với vai trò rõ ràng (system, user, assistant). Hàm mất mát chỉ tính trên phần phản hồi:
\begin{equation}
  \mathcal{L}_{\text{IT}} = -\sum_{t \in \mathcal{R}} \log P(x_t \mid x_{<t}; \theta),
\end{equation}
trong đó $\mathcal{R}$ là tập chỉ số của các token thuộc phần phản hồi. Sinh dữ liệu tổng hợp bằng LLM mạnh hơn để làm dữ liệu instruction tuning cho LLM nhỏ hơn là phương pháp ngày càng phổ biến, đặc biệt trong các miền thiếu dữ liệu chú thích như pháp lý tiếng Việt.

\section{NLP pháp lý tiếng Việt}

NLP pháp lý là lĩnh vực ứng dụng đòi hỏi độ chính xác cao và khả năng suy luận có cấu trúc. Văn bản pháp luật có đặc thù về thuật ngữ chuyên ngành, cấu trúc phân cấp (điều, khoản, điểm) và yêu cầu lập luận chặt chẽ. Đối với tiếng Việt -- một ngôn ngữ ít tài nguyên trong miền pháp lý -- thách thức lớn nhất là sự khan hiếm dữ liệu chú thích chất lượng cao.

Ba tác vụ trong cuộc thi VLSP2025-LegalSML phản ánh các năng lực cốt lõi: (1) \emph{Suy luận tự nhiên} (NLI) yêu cầu xác định điều luật có thể dùng để trả lời câu hỏi hay không; (2) \emph{Trắc nghiệm} (MC) kiểm tra khả năng truy hồi và áp dụng tri thức pháp lý; (3) \emph{Tam đoạn luận} (Syllogism) đòi hỏi mô hình sinh ra lập luận có cấu trúc gồm tiền đề lớn, tiền đề nhỏ và kết luận -- phản ánh đúng bản chất của quá trình áp dụng pháp luật trong thực tế.

\section{Các công trình liên quan}

Trên thế giới, SaulLM-7B~\cite{colombo2024saullm} là mô hình 7B đầu tiên thiết kế riêng cho lĩnh vực luật, được huấn luyện trên 30 tỷ token văn bản pháp lý tiếng Anh. LegalBench~\cite{guha2023legalbench} cung cấp benchmark gồm 162 tác vụ thuộc 6 loại lập luận pháp lý. Đối với tiếng Việt, VLQA~\cite{nguyen2025vlqa} là bộ dữ liệu hỏi đáp pháp lý toàn diện nhất với 3.129 câu hỏi được chuyên gia chú thích. Về tổng hợp dữ liệu, Self-Instruct~\cite{wang2022self} đề xuất sinh ví dụ chỉ dẫn từ mẫu hạt giống, đặt nền móng cho các phương pháp tăng cường dữ liệu dựa trên LLM.

Hai phương pháp đạt thứ hạng cao nhất tại VLSP2025-LegalSML là nguồn cảm hứng trực tiếp cho nghiên cứu. MinLegal~\cite{luan2025minlegal} (hạng 2, score 85,24\%) áp dụng quy trình hai giai đoạn: huấn luyện ba adapter LoRA riêng rồi hợp nhất, sau đó tinh chỉnh toàn phần trên tập gộp; dữ liệu được mở rộng từ 440 lên 3.537 mẫu nhờ sinh few-shot bằng Gemini 2.5 Flash. Bosch@AI~\cite{quang2025bosch} (hạng 1, score 81,08\%) tập trung vào sinh dữ liệu với chuỗi suy luận theo khía cạnh (aspect-based CoT) bằng GPT-4o, mở rộng lên 200.000 ví dụ. Nghiên cứu này kế thừa quy trình hai giai đoạn của MinLegal, đồng thời tích hợp ý tưởng trích xuất khía cạnh của Bosch@AI vào pipeline sinh dữ liệu, hướng tới sự cân bằng giữa chất lượng và chi phí trên môi trường tính toán hạn chế.
